% Technical Abstract -- <=2 pages, SELF-CONTAINED.
% MUST include: name, College, project title (a copy is archived separately).
% Cover: problem, techniques, main results, conclusions.
% Write LAST (5/30).

\section*{Technical Abstract}

\noindent\textbf{Project title:} Directional Reflectance Profile Analysis for Laid-Line Detection in Historical Paper\\
\textbf{Author:} Zhanfeng Zhou\\
\textbf{College:} \emph{TODO}\\
\textbf{Supervisor:} \emph{TODO}\\[0.5em]

% --- The body of the technical abstract: ~2 pages ---
% 1) Problem (paragraph)
%    - Laid-line density in historical paper is a forensic fingerprint.
%    - Manual counting is slow, subjective.
%    - MSI image stacks now exist at scale but no robust automatic pipeline.
%
% 2) Approach (paragraph + small block)
%    - DRP-based 5-stage pipeline (loader, direction, trig-mask,
%      detection, evaluation).
%    - Automatic orientation estimation, patchwise enhancement, multiple
%      detectors (Gabor, multi-phi, spectral-power dual-band).
%
% 3) Results (paragraph with headline numbers)
%    - Synthetic phantom: period error <= +/-2.5\% over 10--50 px,
%      essentially SNR-invariant down to SNR=2.
%    - 9-folio benchmark, two with manual GT: -0.35\% and +0.48\% error.
%    - Failure mode characterised (octave-aliasing at extreme periods).
%
% 4) Conclusions (short paragraph)
%    - Sub-1\% accuracy achievable; pipeline is mould-pair-discriminative
%      in principle.
%    - Future: automatic polarity detection, joint chain-line analysis,
%      mould-mate matching across folios.
